\section{Cost management}
\label{sec:cost}

The cost baseline is the WBS in Figure~\ref{fig:wbs} priced at Level~3. Seventeen control accounts sum to a work estimate of \$1{,}500{,}000. Contingency of \$78{,}890 is sized from residual expected monetary value plus four named residual items, not from a round-number markup, so the authorised budget at completion is \$1{,}578{,}890.

\subsection{Estimating basis and assumptions}
\label{sec:cost-basis}

The estimate is built bottom-up at Level~3. Each package is hours $\times$ an hourly rate, plus plant and materials. Four hourly rates are used. All four start from the same award wage and go through the same four steps.

\begin{enumerate}[leftmargin=1.6em,itemsep=0.25em]
\item \textbf{Wage.} The Manufacturing Award C10 rate is \$29.45 an hour from 1~July~2026 \citep{fwc2026}. It is the legal minimum for a trade-qualified worker. Engineer and specialist rates start from the same wage so all four rates share one published base. C10 is not the award rate for those two jobs.
\item \textbf{Employment on-cost.} The employer pays more than the wage: super 12\%, annual leave 7.69\%, 17.5\% leave loading on that leave, personal leave 3.85\%, public holidays 4.23\%, long service leave 2.17\%, workers compensation and Queensland payroll tax 4.75\%. That adds 37.8\% for office staff and 40.5\% for pad staff, because pad workers compensation is 4.5\% instead of 1.8\%. The multipliers are 1.378 and 1.405.
\item \textbf{Utilisation (assumption).} Not every paid hour can be charged to a package. Briefings, waiting outside the blast zone and paperwork still cost money. We assume 85\% of a technician's paid hours are chargeable, 72\% for engineers and 60\% for specialists. Dividing by that share spreads the unchargeable hours over the chargeable ones.
\item \textbf{Overhead (assumption).} Overhead is the cost of running the crew beyond wages. The multipliers are 1.95 for technicians, 2.66 for engineers, 2.44 for specialists and 1.37 for surge.
\end{enumerate}

\begin{table}[H]
\centering
\small
\caption[Hourly charge-out build-up]{Hourly charge-out build-up. Only the wage and statutory on-costs are published; utilisation and overhead are planning assumptions.}
\label{tab:labour-rates}
\begin{tabular}{@{}lrrrrr@{}}
\toprule
Role & Wage & $\times$ on-cost & $\div$ utilisation & $\times$ overhead & Rate \\
\midrule
Pad technician & 29.45 & 41.38 & 48.68 & 94.93 & \textbf{\$95} \\
Engineer / PM & 29.45 & 40.58 & 56.36 & 149.92 & \textbf{\$150} \\
Specialist & 29.45 & 40.58 & 67.63 & 165.02 & \textbf{\$165} \\
Surge technician$^{\dagger}$ & 44.18 & 62.07 & 73.02 & 100.04 & \textbf{\$100} \\
\bottomrule
\end{tabular}

\vspace{0.2em}
{\footnotesize $^{\dagger}$Surge is paid at overtime first: $\$29.45 \times 1.5 = \$44.18$.}
\end{table}

The crane is priced separately: a 100~t crawler at \$4{,}200 per wet day, operator included, from Queensland plant-hire rates checked against Rawlinsons \citep{rawlinsons2025}.

The budget covers the seventeen packages in Figure~\ref{fig:wbs}. Factory manufacture, engine research and freight are out of scope, so A-100 has no cost. The surge crew (\$22{,}400) is already inside WBS~1.3.3. Option~C's \$21{,}000 is not in the \$1{,}500{,}000; it would come out of the \$78{,}890 reserve only if R-14 happens.

Permits and Juru start on 10~August because they take too long to wait. Other August--September pad work is our choice. By Gate~1 (A-111, 2~October), \$602{,}301 of planned work is already under way: 40.2\% of \$1{,}500{,}000. That is planned work, not cash paid. Gate~1 locks the other \$897{,}699 and the \$78{,}890 contingency rules. After that, adding pad work or moving $T$-0 needs a Board variation.

\subsection{WBS cost baseline}
\label{sec:cost-wbs}

Table~\ref{tab:wbs-cost} is the rollup.
\begin{itemize}[leftmargin=1.4em,itemsep=0.25em]
\item Pad logistics (1.2): \$518{,}000 (34.5\% of the work estimate).
\item Pad testing (1.3): \$414{,}000 (27.6\%).
\item Governance (1.1): \$300{,}000 (20.0\%).
\item Flight/closeout (1.4): \$268{,}000 (17.9\%).
\end{itemize}

\begin{table}[!htbp]
\centering
\footnotesize
\setlength{\tabcolsep}{4.2pt}
\caption[Bottom-up cost baseline by WBS]{Bottom-up cost baseline by Level-3 WBS package. Totals equal the work estimate of \$1{,}500{,}000. A-100 is omitted because it is not a cost account.}
\label{tab:wbs-cost}
\begin{tabularx}{\textwidth}{@{}l p{4.15cm} r X@{}}
\toprule
WBS & Package & Amount (AUD) & Basis \\
\midrule
1.1.1 & PEP and integration & 72{,}000 & 480~h PM @ \$150 \\
1.1.2 & Governance and EVM & 48{,}000 & 320~h planner/cost @ \$150 \\
1.1.3 & ASA / CASA / GBRMPA & 88{,}000 & ASA \$42k $+$ CASA \$28k $+$ GBRMPA \$18k \\
1.1.4 & Juru monitors and heritage & 92{,}000 & Two monitors + survey \\
1.2.1 & Receiving and cleanroom & 68{,}000 & Technician hours + consumables \\
1.2.2 & Crane and 23\,m stacking & 218{,}000 & 8 lift + 16 weld/NDT/bolt-up d $\times$ \$4{,}200 $=$ \$100{,}800 + riggers \\
1.2.3 & Cryo GSE & 180{,}000 & Valves/piping + specialist crew \\
1.2.4 & RF array calibration & 52{,}000 & Crew + analyser hire \\
1.3.1 & Avionics power-on & 58{,}000 & Engineers @ \$150/h + EGSE \\
1.3.2 & WDR (cold) & 160{,}000 & Includes bulk LOX 25~t and LN2 \$68{,}500 \\
1.3.3 & Static fire & 148{,}000 & Includes surge hire \$22{,}400 \\
1.3.4 & Post-fire / pre-kitted spare & 48{,}000 & Engineers + spare handling \\
1.4.1 & LRR and RSO & 38{,}000 & Review board; payload witnesses \\
1.4.2 & Range exclusion / marine & 58{,}000 & Patrol boats + NOTAM \\
1.4.3 & Launch and telemetry capture & 82{,}000 & Launch-day operations \\
1.4.4 & Decrypt and handover & 24{,}000 & Five-day engineering handover \\
1.4.5 & Demobilisation and audit & 66{,}000 & Site restore + audit \\
\midrule
\textbf{Base} & & \textbf{1{,}500{,}000} & \mbox{} \\
\bottomrule
\end{tabularx}

\vspace{0.35em}
{\footnotesize\textit{Notes.} Labour rates are built in Table~\ref{tab:labour-rates}. Crane: 8 lift days and 16 weld/NDT/bolt-up days at the same wet \$4{,}200 (QLD 100~t crawler band about \$3{,}800--\$4{,}500/day), cross-checked against Rawlinsons \citep{rawlinsons2025}. Weather dollars sit only in RES (4 extra wet days). 1.1.3: ASA 280~h @ \$150 $=$ \$42{,}000; CASA CRIS airspace instrument \$28{,}000; GBRMPA permission \$18{,}000 \citep{asa2025,casa2024,gbrmpa2019}. LOX/LN2 \$68{,}500 is a bulk industrial-gas analogue inside 1.3.2. Surge \$22{,}400 $=$ $4 \times 7$~d $\times$ 8~h $\times$ \$100 and sits in 1.3.3, not in contingency.}
\end{table}

The largest package is crane mobilisation and stacking (1.2.2, \$218{,}000) because A-123 is 24 launch-critical days. Table~\ref{tab:wbs-cost} splits the crawler hire from rigger labour. The 16 non-lift days are plant time, not weather. Weather dollars sit only in RES: four extra wet days, \$16{,}800. R-06 is a dawn-lift control with EMV \$0, so it is not a second crane line. The seventeen lines sum to \$1{,}500{,}000. Package 1.2.1 stays one account for A-121 and A-122.

\subsection{Contingency}
\label{sec:cost-contingency}

Contingency is not a flat 15\%. Table~\ref{tab:contingency} adds residual expected monetary value to four named residual items. Treatments already sit in the \$1{,}500{,}000 base (surge hire, dawn-lift window, dual-path recorders, pre-kitted spare), so the pool uses residual $P'_{\%}\times I_{\$}$, not inherent $P_{\%}$. Ten cost-bearing risks sum to \$35{,}760. Itemised RES sums to \$43{,}130. Freight is out of scope and is not in RES. Together they are \$78{,}890 (5.26\% of base), so $BAC$ is \$1{,}578{,}890.

\begin{table}[H]
\centering
\small
\caption[Contingency build-up]{Contingency build-up. The 5.26\% ratio is residual EMV plus four named residual items, not an input markup.}
\label{tab:contingency}
\begin{tabular}{@{}p{8.4cm} r p{4.1cm}@{}}
\toprule
Component & Amount (AUD) & Rule \\
\midrule
$\Sigma$ residual EMV of ten cost-bearing risks & 35{,}760 & Residual $P'_{\%}\times I_{\$}$; treatments already in the base \\
LOX/LN2 price variance (18\% of \$68{,}500 gases) & 12{,}330 & In-scope commodity; not freight \\
Crane weather (only weather \$): 4~d $\times$ \$4{,}200 & 16{,}800 & Extra wet days; R-06 is schedule-only \\
GN2/N2 purge: 40 G-size cylinders $\times$ \$200 & 8{,}000 & Named multiply-out \\
Cryogenic kit: 3 $\times$ 3~m LOX hoses $\times$ \$1{,}400 $+$ seal kit \$1{,}800 & 6{,}000 & Named multiply-out \\
\midrule
\textbf{RES subtotal} & \textbf{43{,}130} & Itemised; not a plug to 15\% \\
\textbf{Contingency reserve} & \textbf{78{,}890} & 5.26\% of base \\
\textbf{$BAC$} & \textbf{1{,}578{,}890} & Base + contingency \\
\bottomrule
\end{tabular}
\end{table}

Several register rows stay out of the \$35{,}760 so the money is not counted twice. R-08 (pad-civil defect) is transferred under the design-and-construct package and carries no Gilmour EMV. R-06, R-11, R-15, R-16, R-17 and R-18 are schedule or HSE exposures with no residual dollar draw. R-12 is the TestFlight~1 in-flight mode; the vehicle is factory product, so site-ops EMV is \$0. If R-14 is realised and Option~C is invoked, the \$21{,}000 is drawn from the \$78{,}890 reserve, not added to the \$1{,}500{,}000 baseline. There is no management reserve above $BAC$: unknown unknowns would need a board variation this plan does not authorise.

\subsection{Time-phased budget and cash flow}
\label{sec:cost-cashflow}

The performance measurement baseline spreads the work estimate on the Table~\ref{tab:precedence} dates. Table~\ref{tab:cashflow} and Figure~\ref{fig:scurve} use a uniform daily burn. October peaks because stacking, Wet Dress Rehearsal and static fire sit there. December is \$6{,}000 because only 1~December of A-145 falls in that month: $\$66{,}000/11 = \$6{,}000$. That is planned closeout, not unused contingency. Contingency is not monthly PV.

\begin{table}[H]
\centering
\small
\caption[Time-phased planned value]{Time-phased planned value of the \$1{,}500{,}000 work baseline. Contingency is not monthly spend.}
\label{tab:cashflow}
\begin{tabular}{@{}lrr@{}}
\toprule
Month & Period PV (AUD) & Cumulative PV (AUD) \\
\midrule
Aug 2026 & 155{,}951 & 155{,}951 \\
Sep 2026 & 407{,}188 & 563{,}139 \\
Oct 2026 (peak: stack, WDR, static fire) & 628{,}861 & 1{,}192{,}000 \\
Nov 2026 & 302{,}000 & 1{,}494{,}000 \\
Dec 2026 & 6{,}000 & \textbf{1{,}500{,}000} \\
\bottomrule
\end{tabular}
\end{table}

Figure~\ref{fig:scurve} ends at \$1{,}500{,}000 and draws $BAC = \$1{,}578{,}890$ as an authorisation line, not December cash. Planned value is \$865{,}578 at M-4 (15~October) and \$602{,}301 at Gate~1. Option~C's \$21{,}000 is outside Table~\ref{tab:cashflow}. If spent, it is actual cost charged to contingency.

\begin{figure}[H]
\centering
\includegraphics[width=0.98\linewidth]{cash_flow_scurve.pdf}
\caption[Planned-value S-curve]{Cumulative planned-value S-curve for the Bowen campaign. Period bars are monthly PV; the teal line is cumulative PV and ends at the \$1{,}500{,}000 work baseline. The crimson dotted line is $BAC$ (\$1{,}578{,}890) as an authorisation ceiling, not December cash. Contingency remains undistributed. Option~C is not in the curve.}
\label{fig:scurve}
\end{figure}
